% ECONOMICS_PROBLEM: {"id": "D91-630", "title": "Biases under high stakes and experience", "jel_code": "D91", "source_id": "OP-0302"}
% !TeX program = xelatex
\documentclass[11pt,a4paper]{article}
\usepackage[margin=23mm]{geometry}
\usepackage{fontspec}
\setmainfont{Latin Modern Roman}
\usepackage{amsmath,amssymb,mathtools}
\usepackage{xcolor,enumitem,booktabs,longtable,array,xurl,hyperref}
\definecolor{muted}{HTML}{53636C}
\hypersetup{colorlinks=true,urlcolor=blue,linkcolor=blue}
\setlength{\parindent}{0pt}
\setlength{\parskip}{5pt}
\newcommand{\R}{\mathbb R}
\newcommand{\E}{\mathbb E}
\newcommand{\N}{\mathbb N}
\newcommand{\ind}{\mathbf 1}
\newcommand{\Var}{\operatorname{Var}}
\newcommand{\Cov}{\operatorname{Cov}}
\newcommand{\supp}{\operatorname{supp}}
\DeclareMathOperator*{\argmin}{arg\,min}
\DeclareMathOperator*{\argmax}{arg\,max}
\newcommand{\entrymeta}[3]{{\sffamily\small\color{muted}\textbf{\texttt{#1}}\quad|\quad #2\par #3\par}}
\newcommand{\Problemref}[1]{\texttt{#1}}
\newcommand{\JELref}[2]{\texttt{#2} (JEL \texttt{#1})}
\begin{document}
\section*{Biases under high stakes and experience}
\textbf{Problem ID:} \texttt{D91-630}\quad \textbf{JEL:} \texttt{D91}\par
% BEGIN ECONOMICS STATEMENT
\label{problem:OP-0302}
{\sffamily\small\color{muted}Original subject group: Behavioral and experimental economics\par}
\label{definitions:problem:30007584}
\textbf{Audit classification:} Background; proposed extension. The genuine 2023 article already tests large incentives in four biases. Professional expertise, repeated feedback and matched market transport are editorial extensions, not an explicitly verified remaining combined agenda.
\subsection*{Definitions and precise research target}
\textbf{Model and research target.} Fix a decision task \(k\) with objectively correct response \(y_k^*(X)\) for problem attributes \(X\). Let \(Y_{ik}(s,t)\) be individual \(i\)'s response when assigned incentive level \(s\geq0\) after \(t\) opportunities to practice with feedback. Define a preregistered bias score \(B_{ik}(s,t)=b_k(Y_{ik}(s,t),y_k^*(X),X)\), where \(b_k\) distinguishes a systematic directional error from unsystematic inaccuracy. Let \(H_i\) denote observed professional expertise and \(P\) a specified target population. The relevant response surface and incentive contrast are
\[\mu_{k,P}(s,t,h)=\mathbb E_P[B_{ik}(s,t)\mid H_i=h],\qquad \tau_{k,P}=\mu_{k,P}(s_H,t,h)-\mu_{k,P}(s_L,t,h).\]
Estimate these quantities using randomized incentives and feedback, adequate support for both stake levels, and payments expressed relative to participants' economic resources. Expertise is not automatically randomized; comparisons across \(h\) require an explicit selection or causal design. Which tasks retain materially positive bias at high stakes after meaningful experience, and how accurately do laboratory estimates predict bias in matched consequential market decisions? Specify thresholds for economic importance and distinguish learning from selective attrition, fatigue, and changes in task interpretation.

\emph{Definitions and maintained assumptions (editorial unless a source definition is named).} A task \(k\) supplies attributes \(X\), an objectively correct answer \(y_k^*(X)\) and a prespecified directional-error score; decisions reflecting legitimate taste differences cannot be assigned a correct answer without a normative model. Normalize incentives by a declared income/resource measure and specify the payment scoring rule: increased nominal maximum prizes do not uniquely define increased marginal incentives. Potential response \(Y_{ik}(s,t)\) refers to assigned stake \(s\) and a controlled training/feedback protocol of length \(t\); experience is an intervention, not elapsed time alone. Expertise \(H_i\) is baseline and may be selected. Define a threshold \(b_0>0\) for material residual bias and contrasts \(\mu(s_H,t,h)-\mu(s_L,t,h)\) only on expertise support. Randomized stakes/feedback identify effects for the sampled population, accounting for repeated measures and attrition; causal expertise effects require another design. Field prediction must preserve task attributes, score meaning and comparable financial consequences. The source’s four task results cannot settle this entire surface.
\subsection*{Variants and assumption changes}
\begin{enumerate}
\item \textbf{Feedback interaction (editorial).} Randomize stakes and the number/type of feedback opportunities factorially. Estimate whether high stakes change learning speed and residual directional bias.
\item \textbf{Expertise transport (editorial).} Compare a prespecified professional population with matched nonprofessionals under identical tasks, and test forecasts on consequential matched decisions. Report predictive differences separately from causal effects of acquiring expertise.
\end{enumerate}
\subsection*{References and source audit}
\begin{enumerate}
\item Introduction p.818; discussion p.830. Supports the stated research area; exact displayed specification is editorial. \url{https://benjamin-enke.com/pdf/Biases_stakes.pdf}
\end{enumerate}
\begin{enumerate}\raggedright
\item\hypertarget{definitions:source:30007584:1}{} Benjamin Enke; Uri Gneezy; Brian Hall; David Martin; Vadim Nelidov; Theo Offerman; Jeroen van de Ven. 2023. \emph{Cognitive Biases: Mistakes or Missing Stakes?}. §I, p. 818; §IV Discussion, p. 830 (PDF pp. 1 and 13)..
\url{https://benjamin-enke.com/pdf/Biases_stakes.pdf}.
DOI: \url{https://doi.org/10.1162/rest_a_01093}.
\end{enumerate}

\subsection*{Common conventions: Source mathematical conventions}

These conventions apply to each entry unless explicitly replaced.

\textbf{Models and identification.} A primitive model \(m\) specifies all probability laws, feasible choices, information, equilibrium and measurement rules used by its target. The admissible class \(\mathcal M\) is fixed independently of outcomes. If \(P_O(m)\) is its observable probability law and \(\tau(m)\) the target, its identified set at observable law \(P\) is
\[
 \mathcal I_\tau(P)=\{\tau(m):m\in\mathcal M,\ P_O(m)=P\}.
\]
Point identification means that this set is a singleton. An empty set signals incompatibility of data and assumptions. Coordinate bounds need not describe the joint set. Bounds are sharp only if every retained target is attainable or arbitrarily approachable by compatible models. Finite bounds require explicit restrictions on unobserved quantities; unrestricted confounding, hidden wealth, extrapolation or arbitrary equilibrium selection can yield uninformative sets.

\textbf{Causal interventions.} A potential outcome \(Y_i(p)\) is the outcome under a complete policy regime \(p\), including timing, financing, information and market spillovers. The target population and units are fixed before treatment. With interference, use the complete assignment vector or a justified exposure mapping; individual no-interference notation alone cannot identify an economy-wide intervention. A difference of mean potential outcomes does not require the cross-world joint distribution. Conditioning on post-treatment migration, survival, enrollment or employment changes the estimand and needs separate justification.

\textbf{Statistical statements.} An estimator, confidence set, forecast or test specifies a sampling law, model class and loss/coverage criterion. Uniform \((1-\alpha)\)-coverage means \(\inf_{m\in\mathcal M}\Pr_m\{\tau(m)\in C_n\}\geq1-\alpha\), or an explicitly asymptotic version. The whole parameter space trivially has coverage, so informativeness requires a stated width, risk or power objective. A forecasting improvement is relative to a named benchmark, information set and evaluation population.

\textbf{Equilibrium and optimization.} An equilibrium comprises feasible plans, optimality, beliefs where needed, budget constraints and all market/resource clearing. The primitives or a stated benchmark define each correspondence \(\mathcal E(p;m)\). Nonemptiness is assumed or proved, never inferred just from compact policy sets. A selected-equilibrium target uses a declared selection rule; a robust target quantifies over all permitted equilibria. Use suprema or infima unless attainment is established. An \(\varepsilon\)-optimal policy is feasible and within \(\varepsilon\) of the finite optimum value. Report an empty feasible set separately.

\textbf{Welfare and accounting.} Normative utility, interpersonal normalization, social weights, numeraire, discounting and the use of public receipts are primitives. Behavioral utility need not coincide with the welfare criterion. Household welfare includes ownership income and rents once; adding profits again to compensating variation generally double-counts them. A monetary equivalent is defined by an explicit transfer and price convention, with monotonicity and range conditions. Average effects, mechanism contributions, transfers and welfare losses are different targets. Infinite sums require convergence and dynamic feasible plans require terminal or no-Ponzi conditions.

\textbf{Source versions.} A working-paper date, revision date and journal publication year are distinguished. A DOI for a journal version is not automatically the DOI of an earlier manuscript. Locators refer to the stated version and printed pages where specified. A metadata-only check does not verify a claim about a full-text conclusion. Every entry's source subsection narrows the provenance claim accordingly.
% END ECONOMICS STATEMENT
\end{document}
